\section{Static and Historical Adaptation Results}

\subsection{Static detector}

The static models use the native LM backend, six-digit precision, raw scoring,
and a one-percent calibration quantile. Their complete development results appear
with the adaptive variants in Table~\ref{tab:master-results}. The main pattern is
high precision with incomplete recall, especially for Chromecast. Accuracy is
reported but should not dominate interpretation because benign windows are much
more common than attack windows.

\begin{table}[H]
\centering
\small
\begin{tabular}{lrrrrr}
\toprule
\textbf{Dataset} & \textbf{$\alpha_0$} & \textbf{$\psi$} &
\textbf{DM cutoff} & \textbf{Dev. windows} & \textbf{Flags} \\
\midrule
D-Link Camera 5 & 122.22 & 0.00818 & $-22.5044$ & 270 & 3 \\
UNSW Chromecast & 28.10 & 0.03559 & $-25.6108$ & 474 & 17 \\
UNSW Samsung camera & 17.39 & 0.05749 & $-16.9221$ & 556 & 52 \\
\bottomrule
\end{tabular}
\caption{Stored static models, composition cutoffs, and development outputs.
Combined events compare their margin score with zero. The three D-Link alerts have no ground-truth classification; they are not
measured true or false positives.}
\label{tab:static-model-summary}
\end{table}

Absolute thresholds are model-specific score values and should not be ranked
across devices. Decisions and labelled metrics provide the valid comparison.

\subsection{Legacy adaptation comparison}

All adaptive numbers in this section belong to the historical legacy rules,
not the current conservative default. The six explicit legacy configurations
and their retained reports preserve these comparisons.
The retained scoring and fixed-reference comparisons are presented separately
in Section~\ref{sec:detector-scoring}; their learning policies and update
schedules must not be conflated with these historical runs.

\begin{figure}[H]
\centering
\begin{tikzpicture}
\begin{axis}[
  ybar,
  bar width=8pt,
  width=0.93\textwidth,
  height=0.36\textwidth,
  ymin=45,
  ymax=85,
  ylabel={F1 score (\%)},
  symbolic x coords={static,threshold,periodic144,periodic288},
  xtick=data,
  xticklabels={Static,Threshold only,Periodic 144/144,Periodic 288/240},
  x tick label style={font=\small,rotate=18,anchor=east},
  y tick label style={font=\small},
  grid=major,
  grid style={gray!20},
  legend style={draw=none,at={(0.5,-0.30)},anchor=north,legend columns=2},
  nodes near coords,
  nodes near coords style={font=\scriptsize},
]
\addplot[fill=linkblue!70] coordinates {
  (static,61.22) (threshold,55.32) (periodic144,58.33) (periodic288,80.70)
};
\addlegendentry{Chromecast}
\addplot[fill=workinggreen!70] coordinates {
  (static,78.69) (threshold,54.00) (periodic144,54.00) (periodic288,77.31)
};
\addlegendentry{Samsung camera}
\end{axis}
\end{tikzpicture}
\caption{Historical static and legacy development comparison. The 288/240 result is selected
development evidence, not a locked final result.}
\label{fig:headline-f1}
\end{figure}


% Confusion counts and metrics are read from the static and explicit legacy JSON reports.
\begin{sidewaystable}[p]
\centering\small
\begin{tabular}{llrrrrrrrrr}
\toprule
\textbf{Device} & \textbf{System} & \textbf{TP} & \textbf{TN} & \textbf{FP} & \textbf{FN} & \textbf{Accuracy} & \textbf{Precision} & \textbf{Recall} & \textbf{F1} & \textbf{FPR} \\
\midrule
Chromecast & Static & 15 & 440 & 2 & 17 & 95.99\% & 88.24\% & 46.88\% & 61.22\% & 0.45\% \\
Chromecast & Legacy threshold, 144 & 13 & 440 & 2 & 19 & 95.57\% & 86.67\% & 40.63\% & 55.32\% & 0.45\% \\
Chromecast & Legacy refit, 144/144 & 14 & 440 & 2 & 18 & 95.78\% & 87.50\% & 43.75\% & 58.33\% & 0.45\% \\
Chromecast & Legacy refit, 288/240 & 23 & 440 & 2 & 9 & 97.68\% & 92.00\% & 71.88\% & 80.70\% & 0.45\% \\
\midrule
Samsung & Static & 48 & 482 & 4 & 22 & 95.32\% & 92.31\% & 68.57\% & 78.69\% & 0.82\% \\
Samsung & Legacy threshold, 144 & 27 & 483 & 3 & 43 & 91.73\% & 90.00\% & 38.57\% & 54.00\% & 0.62\% \\
Samsung & Legacy refit, 144/144 & 27 & 483 & 3 & 43 & 91.73\% & 90.00\% & 38.57\% & 54.00\% & 0.62\% \\
Samsung & Legacy refit, 288/240 & 46 & 483 & 3 & 24 & 95.14\% & 93.88\% & 65.71\% & 77.31\% & 0.62\% \\
\bottomrule
\end{tabular}
\caption{Retained static and legacy development results. FPR is calculated here from the stored FP/TN counts because the older JSON reports predate that output field. The refit labels give update cadence and minimum accepted-window count.}
\label{tab:master-results}
\end{sidewaystable}



Threshold-only adaptation made both devices worse. In the first periodic-refit
configuration, the detector attempted an update every 144 windows and required
144 accepted windows before refitting. It recovered one Chromecast attack
window relative to threshold-only adaptation, while remaining below static;
it made no difference to the Samsung decisions. The model fits
succeeded, so the next experiment tested whether the detector was updating at
the wrong time.

Update intervals from 144 to 432 windows were tested with the same 288-window
buffer and safety margin. Figure~\ref{fig:cadence-sensitivity} shows how F1
changed on both development sets.

\begin{figure}[H]
\centering
\begin{tikzpicture}
\begin{groupplot}[
  group style={group size=2 by 1, horizontal sep=1cm},
  width=0.43\textwidth,
  height=0.32\textwidth,
  xmin=135,
  xmax=440,
  ymin=50,
  ymax=85,
  xtick={144,192,216,240,288,336,360,432},
  x tick label style={font=\scriptsize,rotate=45,anchor=east},
  y tick label style={font=\scriptsize},
  grid=major,
  grid style={gray!25},
  xlabel={Update interval (windows)},
  ylabel={F1 score (\%)},
  title style={font=\bfseries},
  legend style={font=\scriptsize,draw=none,at={(0.5,-0.38)},anchor=north},
]
\nextgroupplot[title={Chromecast}]
\addplot[color=linkblue,mark=*] coordinates {
  (144,55.32) (192,55.32) (216,55.32) (240,61.22)
  (288,66.67) (336,66.67) (360,55.32) (432,55.32)
};
\addlegendentry{Threshold only}
\addplot[color=workinggreen,mark=square*] coordinates {
  (144,58.33) (192,55.32) (216,55.32) (240,61.22)
  (288,80.70) (336,78.57) (360,55.32) (432,55.32)
};
\addlegendentry{Periodic refit}
\addplot[color=black,dashed] coordinates {(144,61.22) (432,61.22)};
\addlegendentry{Static}

\nextgroupplot[title={Samsung camera},ylabel={}]
\addplot[color=linkblue,mark=*] coordinates {
  (144,54.00) (192,60.95) (216,66.06) (240,71.93)
  (288,77.31) (336,79.34) (360,78.69) (432,78.69)
};
\addplot[color=workinggreen,mark=square*] coordinates {
  (144,54.00) (192,59.62) (216,66.06) (240,71.93)
  (288,77.31) (336,79.34) (360,78.69) (432,78.69)
};
\addplot[color=black,dashed] coordinates {(144,78.69) (432,78.69)};
\end{groupplot}
\end{tikzpicture}
\caption{Historical legacy development sensitivity to update cadence, retained from the adaptation-stage report.}
\label{fig:cadence-sensitivity}
\end{figure}


With an update interval of 288 windows, Samsung does not perform the early update
that caused the loss at window 144, while Chromecast still refits once before
the later attacks. The selected configuration requires 240 accepted windows
because both devices had at least that many at the useful update. Figure
\ref{fig:cadence-sensitivity} also shows why this remains a development result:
the apparent best interval depends strongly on when attacks occur in these two
streams.

\subsection{Adaptation safety findings}

After each run, the labels were used to count how many attack windows had been
accepted into the adaptation buffer. The detector itself did not use those
labels when it selected windows or updated its state. This check shows that the
safety gate rejected all Chromecast attacks before the first update, but
accepted 22 Samsung attack windows.

\begin{table}[H]
\centering
\small
\begin{tabular}{llrrr}
\toprule
\textbf{Device} & \textbf{Update} & \textbf{Accepted windows} &
\textbf{Attack windows} & \textbf{Contamination} \\
\midrule
Chromecast & 288 & 285 & 0 & 0.00\% \\
Chromecast & 432 & 288 & 3 & 1.04\% \\
Samsung camera & 288 & 241 & 22 & 9.13\% \\
Samsung camera & 432 & 288 & 20 & 6.94\% \\
\bottomrule
\end{tabular}
\caption{Attack windows found in each periodic-refit buffer after the run.}
\label{tab:safe-contamination}
\end{table}

Increasing the score margin alone did not remove the Samsung contamination.
Those attacks already look ordinary to a detector that observes only four
protocol proportions. This supports an independent volume or rate safeguard
rather than a larger value of the same composition-score margin.

Refitting also changes the total concentration $\alpha_0$. The Chromecast value
moves moderately, whereas the Samsung value becomes more than nine times larger
under the selected configuration.

\begin{table}[H]
\centering
\small
\begin{tabularx}{\textwidth}{l>{\raggedright\arraybackslash}Xrrr}
\toprule
\textbf{Device} & \textbf{Experiment} & \textbf{Initial $\alpha_0$} &
\textbf{Final $\alpha_0$} & \textbf{Ratio} \\
\midrule
Chromecast & Periodic refit, 144/144 & 28.10 & 36.05 & 1.28 \\
Chromecast & Periodic refit, 288/240 (selected) & 28.10 & 34.47 & 1.23 \\
Samsung camera & Periodic refit, 144/144 & 17.39 & 206.40 & 11.87 \\
Samsung camera & Periodic refit, 288/240 (selected) & 17.39 & 166.62 & 9.58 \\
\bottomrule
\end{tabularx}
\caption{Movement of Dirichlet concentration during promoted refits. In each
periodic-refit name, the first number is the update interval and the second is
the minimum number of accepted windows required for refitting.}
\label{tab:alpha-movement}
\end{table}

Convergence and a better fit to the accepted buffer are therefore not enough to
make a refitted model safe. These findings motivated the concentration, proportion, threshold, and volume
guards now implemented by the conservative policy. They do not measure how
often the new policy will learn successfully on representative traffic.

\subsection{LM backend evidence}

Figure~\ref{fig:lm-runtime} compares LM and SciPy directly on the D-Link,
Chromecast, and Samsung training matrices. This three-dataset comparison is retained from the report of
9 September 2026 and is not a new timing run. Only the D-Link production
benchmark JSON is retained at its current reference path; its exact stored
timings are given separately below. Model fitting is
shown in milliseconds; the two faster operations are shown in microseconds.

\begin{figure}[H]
\centering
\begin{tikzpicture}
\begin{groupplot}[
  group style={group size=3 by 1, horizontal sep=0.8cm},
  ybar,
  /pgf/bar width=5pt,
  width=0.29\textwidth,
  height=0.31\textwidth,
  ymin=0,
  symbolic x coords={dlink,chromecast,samsung},
  xtick=data,
  xticklabels={D-Link,Chromecast,Samsung},
  x tick label style={font=\scriptsize,rotate=28,anchor=east},
  y tick label style={font=\scriptsize},
  grid=major,
  grid style={gray!20},
  enlarge x limits=0.22,
  title style={font=\small\bfseries},
  legend style={font=\small,draw=none,at={(1.72,-0.36)},anchor=north,
    legend columns=2},
]
\nextgroupplot[title={Complete model fit},ylabel={Time (ms)},ymax=285]
\addplot[fill=linkblue!70] coordinates {
  (dlink,191.04) (chromecast,254.17) (samsung,42.46)
};
\addlegendentry{LM}
\addplot[fill=workinggreen!70] coordinates {
  (dlink,180.59) (chromecast,194.40) (samsung,34.47)
};
\addlegendentry{SciPy}

\nextgroupplot[title={Training likelihood},ylabel={Time ($\mu$s)},ymax=95]
\addplot[fill=linkblue!70] coordinates {
  (dlink,55.02) (chromecast,84.69) (samsung,45.69)
};
\addplot[fill=workinggreen!70] coordinates {
  (dlink,49.79) (chromecast,52.12) (samsung,31.57)
};

\nextgroupplot[title={Single-window score},ylabel={Time ($\mu$s)},ymax=15]
\addplot[fill=linkblue!70] coordinates {
  (dlink,12.00) (chromecast,12.22) (samsung,12.18)
};
\addplot[fill=workinggreen!70] coordinates {
  (dlink,13.18) (chromecast,13.22) (samsung,13.20)
};
\end{groupplot}
\end{tikzpicture}
\caption{Historical median LM and SciPy runtimes for three production workloads. A complete
fit excludes data loading and file output; training likelihood is one call on
the complete fit matrix; single-window scoring uses one non-silent window. Lower
is better.}
\label{fig:lm-runtime}
\end{figure}


\begin{table}[H]
\centering
\small
\begin{tabular}{lr}
\toprule
\textbf{Agreement measure} & \textbf{Stored value} \\
\midrule
Maximum relative alpha difference & $3.86\times10^{-7}$ \\
Maximum absolute score difference & $1.09\times10^{-5}$ \\
Decision mismatches & 0 of 270 \\
Flags from each backend & 3 \\
\bottomrule
\end{tabular}
\caption{LM and SciPy agreement in the D-Link production benchmark.}
\label{tab:lm-agreement}
\end{table}

For the retained combined detector, the separate
\path{reports/detection_backend_checks.json} record also compares LM and SciPy
on all three configured datasets. It reports zero decision mismatches across
their 1,300 development windows; that agreement does not itself measure an
improvement in detection accuracy.

The native matrix optimisation removes repeated LM setup while preserving the
row-wise model. A separate experiment reproduces the aggregate calculation used
by the reference study.

\begin{table}[H]
\centering
\small
\begin{tabularx}{\textwidth}{Xrrr}
\toprule
\textbf{Architecture} & \textbf{\shortstack{Paired fits\\converged}} &
\textbf{\shortstack{Likelihood-time\\reduction}} &
\textbf{\shortstack{Wall-time\\reduction}} \\
\midrule
Reference-paper aggregate & 79 of 84 & 12.2\% & 8.8\% \\
\bottomrule
\end{tabularx}
\caption{LM timing result for the aggregate reference-paper experiment.}
\label{tab:professor-benchmark}
\end{table}

That result applies only to the aggregate calculation. In the production
benchmark, LM is slower for complete fitting and matrix likelihood on all three
datasets, but slightly faster for single-window scoring. The numerical agreement
in Table~\ref{tab:lm-agreement} shows that the backend choice does not change the
D-Link decisions.
